%% secoes/00b-resumo.tex — incluído por main.tex
%%
%% A SBC pede, para artigos em português, o abstract (em inglês) e o
%% resumo (em português), cada um com no máximo 10 linhas, ambos na
%% primeira página. O template SBC não tem campo de palavras-chave.

\begin{abstract}
  This study identifies characteristics associated with low birth weight, preterm birth, congenital anomaly and low 5-minute Apgar score in vulnerable Brazilian populations. The SSDP+ subgroup discovery algorithm was applied to more than 5 million births from 2011 to 2015 in the 100 Million Brazilian Cohort linked to SINASC, in the overall sample and seven vulnerable segments (32 analyses). Few prenatal visits were the most frequent risk factor for low birth weight, preterm birth and low Apgar score, reaching, alone, 2.03 times the baseline, and first pregnancy accompanied these three outcomes. Congenital anomaly showed a distinct pattern, marked by cesarean delivery, male sex and better household conditions. Mothers under 15 had the highest prevalences of low birth weight, preterm birth and anomaly, and home birth stood out among Indigenous births.
\end{abstract}

\begin{resumo}
  Este trabalho identifica características associadas ao baixo peso ao nascer, ao nascimento pré-termo, à malformação congênita e ao Apgar baixo no quinto minuto em populações vulneráveis do Brasil. Aplicou-se o algoritmo SSDP+ de descoberta de subgrupos a mais de 5 milhões de nascimentos de 2011 a 2015 da Coorte de 100 Milhões de Brasileiros vinculada ao SINASC, no recorte geral e em sete segmentos vulneráveis (32 análises). Poucas consultas de pré-natal foram o fator mais recorrente no baixo peso, no pré-termo e no Apgar baixo, chegando, isoladamente, a 2,03 vezes a basal, e a primeira gestação acompanhou esses três desfechos. A malformação teve padrão distinto, marcado por parto cesáreo, sexo masculino e melhores condições domiciliares. Mães menores de 15 anos tiveram as maiores prevalências de baixo peso, pré-termo e malformação, e o nascimento domiciliar se destacou entre indígenas.
\end{resumo}
